\SourceLeaf{171}
\begin{center}
\footnotesize\setlength{\tabcolsep}{2pt}
\begin{tabular}{r|l|r|r||r|l|r|r||r|l|r|r}
135&O.&102 Thl.&1179 Thl.&170&R.&124 Thl.&1261 Thl.&199&F.&167 Thl.&937 Thl.\\
138&C.&82&1402&171&C.&132&955&200&L.&211&748\\
139&U.&137&730&173&F.&168&931&201&W.&194&749\\
140&H.&134&1199&174&H.&163&1354&202&M.&201&991\\
141&L.&113&833&175&R.&195&996&203&B.&167&751\\
147&R.&133&1188&176&L.&190&1122&204&E.&201&697\\
149&G.&127&1439&177&W.&176&936&205&H.&195&790\\
151&O.&173&870&180&T.&177&887&206&H.&188&854\\
153&B.&154&1271&185&W.&198&919&207&W.&200&783\\
155&K.&145&727&186&H.&169&1134&208&E.&175&1278\\
156&v. S.&146&662&187&W.&180&748&209&B.&154&1051\\
160&H.&161&1173&189&H.&172&953&211&S.&195&1164\\
161&B.&156&682&192&R.&188&782&212&T.&181&895\\
164&W.&159&1155&193&G.&206&896&213&W.&200&811\\
167&S.&190&835&195&D.&191&822&214&B.&188&915\\
169&Z.&180&943&196&B.&207&968&&&&
\end{tabular}
\end{center}
\noindent\hspace*{\parindent}Zum Verständniss der Positionen (3) und (8) dient Folgendes. Das älteste der jetzigen Mitglieder hat die Numerirung 135: die ersten 134 sind sämmtlich verstorben oder auf andere Weise ausgeschieden. Für jedes dieser 134 ausgeschiedenen Mitglieder ist berechnet: der Betrag der von demselben geleisteten Beiträge, und, wo Hinterbliebene Pensionen erhalten haben, der Betrag der Witwenpensionen und der Waisenpensionen. Alle diese Zahlungen sind aber nicht nach ihrer wirklichen Grösse in Ansatz gebracht, sondern nach den gegenwärtig bestehenden Sätzen, nemlich 15 Thaler für jährlichen Beitrag, und 200 Thaler als jetziger Betrag der festen Pension: ausserdem aber sind sämmtliche gezahlten oder empfangenen Gelder durch Discontirung auf das Zeitmoment des Eintritts des betreffenden Mitgliedes reducirt. Rücksichtlich der Modificationen, welche an diese Rechnungen noch angebracht werden mussten bei denjenigen, aus diesen 134 Mitgliedern, von welchen die Witwen noch am Leben sind, beziehe ich mich auf meine Denkschrift von 1845.

So verstanden, beträgt für alle diese 134 Mitglieder
\begin{center}
\begin{tabular}{l r l}
die Summe aller Beiträge . . . . . .&25506&Thaler\\
die Summe aller Witwenpensionen&66771&,,\\
die Summe aller Waisenpensionen&9604&,,
\end{tabular}
\end{center}
\noindent woraus sich die Durchschnittswerthe ergeben: für die Beiträge 190,35 Thaler, für die Witwenpensionen 498,29 Thaler, für die Waisenpensionen 71,67 Thaler. Im Jahr 1845, wo die bis dahin ausgeschiedenen Mitglieder sich in ununterbrochener Folge nur bis zu Nr. 108 erstreckten, hatte ich für diese eine ganz ähnliche Rechnung ausgeführt, welche für die Durchschnittswerthe, wenn sie in dieselbe Form gebracht werden, ergeben hatte: 193,18 Thaler; 503,08 Thaler; 81,07 Thaler. Man erkennt hieraus mit Befriedigung, dass die bisherigen Erfahrungen sich schon sehr wohl zur Feststellung von Durchschnittswerthen eignen, und darin der Abschätzung der künftigen Bedürfnisse eine werthvolle Grundlage geben. Nach dem neuen Resultate ist das Verhältniss der Waisenpension zur Witwenpension durchschnittlich sehr nahe \(\tfrac{1}{7}\), welches zur Feststellung der Position (2) benutzt ist (man vergl. S. [165] am Schluss).

Da nun ferner den allmähligen Leistungen jedes neu eintretenden Mitgliedes einerseits, und den Pensionsbezügen durch seine Relicten andererseits, die Summen 190,35 und 569,96 Thaler, im Zeitpunkte des Eintritts einmal in die Kasse eingezahlt und resp. aus derselben ausgezahlt, nach den Durchschnittswerthen äquivaliren: so brauchte man, wenn jedes Jahr Ein neues Mitglied einträte, nur diese beiden Zahlungen jedes Jahr wiederholt zu denken, und sie nun zu capitalisiren, d. i. ein für allemal der Kasse theils die Capitalsumme 5439 Thaler als Einnahme zuzuweisen, theils die andern 16285 Thaler als Ausgabe anzurechnen, um den Verpflichtungen und Berechtigungen aller künftig beitretenden Mitglieder Rechnung zu tragen. Diese Summen müssen nun aber noch mit derjenigen Ziffer multiplicirt werden, welche die Durch%
